% GENERATED FILE. Edit canonical lessons, then run npm run generate:books.
% canonical-source-hash: fnv1a64:82410e010f914649
% canonical-lessons: KA-C361-smaraka, KA-C361-samadhi, KA-C361-smarisu, KA-C361-buruju, KA-C361-hulu, KA-R361-review-355-361

\chapter{Remembering the Fallen}
\label{ch:ka-remembering-the-fallen}
\hlchaptermodality{361}

\begin{chapteropening}
\textbf{By the end of this chapter:} \emph{I can say a memorial, a tomb, to commemorate, a bastion and to bury.}

Complete the last lesson of chapter 361: I can say a memorial, a tomb, to commemorate, a bastion and to bury.
\end{chapteropening}

\section[\texorpdfstring{\kn{ಸ್ಮಾರಕ} (smāraka)}{ಸ್ಮಾರಕ (smāraka)}]{\kn{ಸ್ಮಾರಕ} (smāraka) — a memorial; a monument}

\label{lesson:KA-C361-smaraka}



\begin{quote}
Before the new one: say the Kannada for to barge in, then the Kannada for to snap.
\end{quote}

\subsection*{You'll want to know: \kn{ಸ್ಮಾರಕ}}
\textbf{\kn{ಸ್ಮಾರಕ}} — \emph{smāraka} — \textquotedblleft{}a memorial; a monument\textquotedblright{}.

\begin{quote}
\kn{ಊರಿನ} \kn{ಮಧ್ಯದಲ್ಲಿ} \kn{ಸೈನಿಕರ} \kn{ಸ್ಮಾರಕ} \kn{ಇದೆ}.
\end{quote}

\begin{quote}
\emph{ūrina madhyadalli sainikara smāraka ide.}
\end{quote}

\begin{quote}
There is a soldiers' memorial in the middle of the town.
\end{quote}

Plural: \kn{ಸ್ಮಾರಕಗಳು}, \emph{smārakagaḷu}.

Old temples and forts are \kn{ಐತಿಹಾಸಿಕ} \kn{ಸ್ಮಾರಕಗಳು}, historic monuments.

\begin{cousinweb}
\textbf{\kn{ಸ್ಮಾರಕ}} (\emph{smāraka}) is borrowed from Sanskrit \emph{smāraka}, \textquotedblleft{}recalling, reminding of\textquotedblright{} (Monier-Williams), from \emph{smṛ}, \textquotedblleft{}to remember\textquotedblright{}.
\end{cousinweb}

\begin{grammarlens}[title={What you've built}]
One more word for telling what happened.
\end{grammarlens}

\subsection*{Your turn}
\begin{itemize}
\raggedright
  \item \emph{Say it:} \emph{smāraka}
  \item \emph{Say it:} \emph{smāraka}, once more
  \item \emph{Say it:} \emph{ūrina madhyadalli sainikara smāraka ide.} — There is a soldiers' memorial in the middle of the town.
  \item \emph{Recall:} say the Kannada for to barge in, then the Kannada for to snap, then say \emph{smāraka} again
\end{itemize}

\begin{tcolorbox}[breakable,colback=teal!4,colframe=teal!35!black,title={Before you move on}]
What does \kn{ಸ್ಮಾರಕ} mean? (\textquotedblleft{}A memorial; a monument\textquotedblright{}.) Say it once more.
\end{tcolorbox}
\section[\texorpdfstring{\kn{ಸಮಾಧಿ} (samādhi)}{ಸಮಾಧಿ (samādhi)}]{\kn{ಸಮಾಧಿ} (samādhi) — a tomb; a grave}

\label{lesson:KA-C361-samadhi}



\begin{quote}
Before the new one: say the Kannada for to snap, then the Kannada for a memorial.
\end{quote}

\subsection*{You'll want to know: \kn{ಸಮಾಧಿ}}
\textbf{\kn{ಸಮಾಧಿ}} — \emph{samādhi} — \textquotedblleft{}a tomb; a grave\textquotedblright{}.

\begin{quote}
\kn{ಹಳೆಯ} \kn{ರಾಜನ} \kn{ಸಮಾಧಿ} \kn{ಬೆಟ್ಟದ} \kn{ಕೆಳಗೆ} \kn{ಇದೆ}.
\end{quote}

\begin{quote}
\emph{haḷeya rājana samādhi beṭṭada keḷage ide.}
\end{quote}

\begin{quote}
The old king's tomb is below the hill.
\end{quote}

Plural: \kn{ಸಮಾಧಿಗಳು}, \emph{samādhigaḷu}.

In religion \kn{ಸಮಾಧಿ} is also deep meditation.

\begin{cousinweb}
\textbf{\kn{ಸಮಾಧಿ}} (\emph{samādhi}) is borrowed from Sanskrit \emph{samādhi}, \textquotedblleft{}putting together, joining\textquotedblright{}, with later senses that include deep meditation and a tomb (Monier-Williams).
\end{cousinweb}

\begin{grammarlens}[title={What you've built}]
One more word for telling what happened.
\end{grammarlens}

\subsection*{Your turn}
\begin{itemize}
\raggedright
  \item \emph{Say it:} \emph{samādhi}
  \item \emph{Say it:} \emph{samādhi}, once more
  \item \emph{Say it:} \emph{haḷeya rājana samādhi beṭṭada keḷage ide.} — The old king's tomb is below the hill.
  \item \emph{Recall:} say the Kannada for to snap, then the Kannada for a memorial, then say \emph{samādhi} again
\end{itemize}

\begin{tcolorbox}[breakable,colback=teal!4,colframe=teal!35!black,title={Before you move on}]
What does \kn{ಸಮಾಧಿ} mean? (\textquotedblleft{}A tomb; a grave\textquotedblright{}.) Say it once more.
\end{tcolorbox}
\section[\texorpdfstring{\kn{ಸ್ಮರಿಸು} (smarisu)}{ಸ್ಮರಿಸು (smarisu)}]{\kn{ಸ್ಮರಿಸು} (smarisu) — to commemorate; to remember}

\label{lesson:KA-C361-smarisu}



\begin{quote}
Before the new one: say the Kannada for a memorial, then the Kannada for a tomb.
\end{quote}

\subsection*{You'll want to know: \kn{ಸ್ಮರಿಸು}}
\textbf{\kn{ಸ್ಮರಿಸು}} — \emph{smarisu} — \textquotedblleft{}to commemorate; to remember\textquotedblright{}.

\begin{quote}
\kn{ಪ್ರತಿ} \kn{ವರ್ಷ} \kn{ಊರಿನವರು} \kn{ಆ} \kn{ಯೋಧರನ್ನು} \kn{ಸ್ಮರಿಸುತ್ತಾರೆ}.
\end{quote}

\begin{quote}
\emph{prati varṣa ūrinavaru ā yōdharannu smarisuttāre.}
\end{quote}

\begin{quote}
Every year the townspeople commemorate those warriors.
\end{quote}

Past (I …): \kn{ಸ್ಮರಿಸಿದೆ}, \emph{smariside}.

Who is remembered takes -\kn{ಅನ್ನು}. A formal word; in talk, \kn{ನೆನಪಿಸಿಕೊಳ್ಳು}.

\begin{cousinweb}
\textbf{\kn{ಸ್ಮರಿಸು}} (\emph{smarisu}) is built on Sanskrit \emph{smaraṇa}, \textquotedblleft{}remembering, memory\textquotedblright{} (CDIAL 13862), with the Kannada verb ending -\emph{isu}.
\end{cousinweb}

\begin{grammarlens}[title={What you've built}]
One more word for telling what happened.
\end{grammarlens}

\subsection*{Your turn}
\begin{itemize}
\raggedright
  \item \emph{Say it:} \emph{smarisu}
  \item \emph{Say it:} \emph{smarisu}, once more
  \item \emph{Say it:} \emph{prati varṣa ūrinavaru ā yōdharannu smarisuttāre.} — Every year the townspeople commemorate those warriors.
  \item \emph{Recall:} say the Kannada for a memorial, then the Kannada for a tomb, then say \emph{smarisu} again
\end{itemize}

\begin{tcolorbox}[breakable,colback=teal!4,colframe=teal!35!black,title={Before you move on}]
What does \kn{ಸ್ಮರಿಸು} mean? (\textquotedblleft{}To commemorate; to remember\textquotedblright{}.) Say it once more.
\end{tcolorbox}
\section[\texorpdfstring{\kn{ಬುರುಜು} (buruju)}{ಬುರುಜು (buruju)}]{\kn{ಬುರುಜು} (buruju) — a bastion; a tower}

\label{lesson:KA-C361-buruju}



\begin{quote}
Before the new one: say the Kannada for a tomb, then the Kannada for to commemorate.
\end{quote}

\subsection*{You'll want to know: \kn{ಬುರುಜು}}
\textbf{\kn{ಬುರುಜು}} — \emph{buruju} — \textquotedblleft{}a bastion; a tower\textquotedblright{}.

\begin{quote}
\kn{ಕೋಟೆಯ} \kn{ಬುರುಜಿನ} \kn{ಮೇಲೆ} \kn{ಸೈನಿಕರ} \kn{ಸ್ಮಾರಕ} \kn{ಇದೆ}.
\end{quote}

\begin{quote}
\emph{kōṭeya burujina mēle sainikara smāraka ide.}
\end{quote}

\begin{quote}
There is a soldiers' memorial on the fort's bastion.
\end{quote}

Plural: \kn{ಬುರುಜುಗಳು}, \emph{burujugaḷu}.

The round towers on a fort wall where guns stood.

\begin{cousinweb}
\textbf{\kn{ಬುರುಜು}} (\emph{buruju}) is a loan from Arabic \emph{burj}, \textquotedblleft{}a tower\textquotedblright{}, which came through Persian and Hindustani.
\end{cousinweb}

\begin{grammarlens}[title={What you've built}]
One more word for telling what happened.
\end{grammarlens}

\subsection*{Your turn}
\begin{itemize}
\raggedright
  \item \emph{Say it:} \emph{buruju}
  \item \emph{Say it:} \emph{buruju}, once more
  \item \emph{Say it:} \emph{kōṭeya burujina mēle sainikara smāraka ide.} — There is a soldiers' memorial on the fort's bastion.
  \item \emph{Recall:} say the Kannada for a tomb, then the Kannada for to commemorate, then say \emph{buruju} again
\end{itemize}

\begin{tcolorbox}[breakable,colback=teal!4,colframe=teal!35!black,title={Before you move on}]
What does \kn{ಬುರುಜು} mean? (\textquotedblleft{}A bastion; a tower\textquotedblright{}.) Say it once more.
\end{tcolorbox}
\section[\texorpdfstring{\kn{ಹೂಳು} (hūḷu)}{ಹೂಳು (hūḷu)}]{\kn{ಹೂಳು} (hūḷu) — to bury}

\label{lesson:KA-C361-hulu}



\begin{quote}
Before the new one: say the Kannada for to commemorate, then the Kannada for a bastion.
\end{quote}

\subsection*{You'll want to know: \kn{ಹೂಳು}}
\textbf{\kn{ಹೂಳು}} — \emph{hūḷu} — \textquotedblleft{}to bury\textquotedblright{}.

\begin{quote}
\kn{ಅವರು} \kn{ಸತ್ತ} \kn{ನಾಯಿಯನ್ನು} \kn{ತೋಟದಲ್ಲಿ} \kn{ಹೂಳಿದರು}.
\end{quote}

\begin{quote}
\emph{avaru satta nāyiyannu tōṭadalli hūḷidaru.}
\end{quote}

\begin{quote}
They buried the dead dog in the garden.
\end{quote}

Past (I …): \kn{ಹೂಳಿದೆ}, \emph{hūḷide}.

What is buried takes -\kn{ಅನ್ನು}; where takes -\kn{ಅಲ್ಲಿ}.

\begin{cousinweb}
\textbf{\kn{ಹೂಳು}} (\emph{hūḷu}) is native Dravidian: DEDR 4376 lists the older Kannada forms of this verb, \textquotedblleft{}to cover, to bury\textquotedblright{}, beginning with \emph{p-} and later \emph{h-}.
\end{cousinweb}

\begin{grammarlens}[title={What you've built}]
One more word for telling what happened.
\end{grammarlens}

\subsection*{Your turn}
\begin{itemize}
\raggedright
  \item \emph{Say it:} \emph{hūḷu}
  \item \emph{Say it:} \emph{hūḷu}, once more
  \item \emph{Say it:} \emph{avaru satta nāyiyannu tōṭadalli hūḷidaru.} — They buried the dead dog in the garden.
  \item \emph{Recall:} say the Kannada for to commemorate, then the Kannada for a bastion, then say \emph{hūḷu} again
\end{itemize}

\begin{tcolorbox}[breakable,colback=teal!4,colframe=teal!35!black,title={Before you move on}]
What does \kn{ಹೂಳು} mean? (\textquotedblleft{}To bury\textquotedblright{}.) Say it once more.
\end{tcolorbox}
\section[(dialogue)]{Looking back: from Digging Up the Past to Remembering the Fallen}

\label{lesson:KA-R361-review-355-361}



\begin{quote}
Say the Kannada for an excavation and archaeology.
\end{quote}

\subsection*{Your turn}
Say these aloud:

\begin{itemize}
\raggedright
  \item an excavation, archaeology, ruins, an inscription, a manuscript
  \item folklore, a vassal king, a battle, a ballad, a moral
  \item to swindle, confidential, to cover up, to reveal, a plot
  \item a hearing, a lawsuit, a case, bail, to confess
  \item a public festival, a procession, a slogan, to parade, a flag
  \item to hunt, to wander, to lie in wait, to barge in, to snap
  \item a memorial, a tomb, to commemorate, a bastion, to bury
\end{itemize}

\begin{tcolorbox}[breakable,colback=teal!4,colframe=teal!35!black,title={Before you move on}]
An excavation? (\textbf{\kn{ಉತ್ಖನನ}}, \emph{utkhanana}.) Archaeology? (\textbf{\kn{ಪುರಾತತ್ವ}}, \emph{purātatva}.) Ruins? (\textbf{\kn{ಅವಶೇಷ}}, \emph{avaśēṣa}.) An inscription? (\textbf{\kn{ಶಾಸನ}}, \emph{śāsana}.) A manuscript? (\textbf{\kn{ಹಸ್ತಪ್ರತಿ}}, \emph{hastaprati}.) Folklore? (\textbf{\kn{ಜಾನಪದ}}, \emph{jānapada}.) A vassal king? (\textbf{\kn{ಸಾಮಂತ}}, \emph{sāmanta}.) A battle? (\textbf{\kn{ಸಮರ}}, \emph{samara}.) A ballad? (\textbf{\kn{ಲಾವಣಿ}}, \emph{lāvaṇi}.) A moral? (\textbf{\kn{ನೀತಿ}}, \emph{nīti}.) To swindle? (\textbf{\kn{ವಂಚಿಸು}}, \emph{vañcisu}.) Confidential? (\textbf{\kn{ಗೌಪ್ಯ}}, \emph{gaupya}.) To cover up? (\textbf{\kn{ಮುಚ್ಚಿಡು}}, \emph{mucciḍu}.) To reveal? (\textbf{\kn{ಬಹಿರಂಗಪಡಿಸು}}, \emph{bahiraṅgapaḍisu}.) A plot? (\textbf{\kn{ಸಂಚು}}, \emph{sañcu}.) A hearing? (\textbf{\kn{ವಿಚಾರಣೆ}}, \emph{vicāraṇe}.) A lawsuit? (\textbf{\kn{ಮೊಕದ್ದಮೆ}}, \emph{mokaddame}.) A case? (\textbf{\kn{ಪ್ರಕರಣ}}, \emph{prakaraṇa}.) Bail? (\textbf{\kn{ಜಾಮೀನು}}, \emph{jāmīnu}.) To confess? (\textbf{\kn{ತಪ್ಪೊಪ್ಪಿಕೊ}}, \emph{tappoppiko}.) A public festival? (\textbf{\kn{ಉತ್ಸವ}}, \emph{utsava}.) A procession? (\textbf{\kn{ಮೆರವಣಿಗೆ}}, \emph{meravaṇige}.) A slogan? (\textbf{\kn{ಘೋಷವಾಕ್ಯ}}, \emph{ghōṣavākya}.) To parade? (\textbf{\kn{ಮೆರೆಸು}}, \emph{meresu}.) A flag? (\textbf{\kn{ಧ್ವಜ}}, \emph{dhvaja}.) To hunt? (\textbf{\kn{ಬೇಟೆಯಾಡು}}, \emph{bēṭeyāḍu}.) To wander? (\textbf{\kn{ಅಲೆದಾಡು}}, \emph{aledāḍu}.) To lie in wait? (\textbf{\kn{ಹೊಂಚು} \kn{ಹಾಕು}}, \emph{hoñcu hāku}.) To barge in? (\textbf{\kn{ನುಗ್ಗು}}, \emph{nuggu}.) To snap? (\textbf{\kn{ಮುರಿ}}, \emph{muri}.) A memorial? (\textbf{\kn{ಸ್ಮಾರಕ}}, \emph{smāraka}.) A tomb? (\textbf{\kn{ಸಮಾಧಿ}}, \emph{samādhi}.) To commemorate? (\textbf{\kn{ಸ್ಮರಿಸು}}, \emph{smarisu}.) A bastion? (\textbf{\kn{ಬುರುಜು}}, \emph{buruju}.) To bury? (\textbf{\kn{ಹೂಳು}}, \emph{hūḷu}.)
\end{tcolorbox}
